

\section{模型评价与推广}

\subsection{模型优点}

本文围绕三通道数据组织计算，方法选择与观测条件相适应。伪迹筛选将协方差结构与分频段异常结合，对位置提示和形状提示任务采用不同检测器配置；阈值敏感性和等样本量抽样同时检查清洗强度及样本组成变化，便于发现局部峰值受筛选影响的情况。P-spline 则给出完整 ERP 轨迹，保留了候选峰以外的响应信息。

正向模型保留边缘方向及空间位置，可逐级追踪视觉输入经群体活动形成电极差异的过程，并以空间消融检查模型内部的差异来源。认知模型直接从实测脑电构造功能代理，便于对照各频段的贡献。评价中分别报告总体波形、左右差异波和记录留出分类结果，避免仅凭平均曲线吻合就推断单试次判别有效。

\subsection{模型不足}

三电极对五源的映射不唯一，固定源坐标和观测模式变换限制了生理解释。总体 ERP 虽取得较高样本内相关性，左右差异波仍未复现；六组参数拟合仅一组报告收敛，部分候选参数接近搜索边界。第一问的局部峰值对筛选阈值敏感，较强清洗可能改变条件样本组成。第三问代理降低了特征重建误差，但分类及应答侧预测未呈现稳定增益；其重建结果也不能作为独立脑源的证据。候选目标时刻、有限记录数及尚待核实的行为标记进一步限制了结论范围。

\subsection{模型推广}

该流程可作为同类少通道视觉实验的分析起点：更换任务时，重新确定事件窗、检测频段与阈值，并检查保留样本的构成。正向模型可通过更换形状模板研究其他视觉图形，再以独立记录检验所预测的差异。若补充顶枕区电极、个体头模型及可靠目标触发标记，可进一步约束映射并检验认知代理与行为的关系。用于脑机接口或疾病研究前，还需扩大受试者样本、加入相应标签并开展独立验证，现有结果尚不足以支持临床诊断。
